\section{Especialización y generalización}

El modelo utiliza una especialización de \texttt{Empleado} en
\texttt{Gerente}, \texttt{Cajero}, \texttt{Encargado de premios} y
\texttt{Técnico}. El símbolo \texttt{d} y la línea doble entre la
superentidad y el triángulo establecen dos restricciones:

\begin{itemize}[leftmargin=*]
  \item \textbf{disjunta}: un empleado sólo puede pertenecer a una de las
        cuatro subentidades;
  \item \textbf{total}: todo empleado registrado debe desempeñar exactamente
        uno de esos puestos.
\end{itemize}

Las subentidades heredan \texttt{CURP}, \texttt{RFC}, nombre completo, fecha
de nacimiento, teléfonos, correos electrónicos, fecha de contratación y
salario. Además, cada una conserva los atributos o relaciones propios de su
función:

\begin{itemize}[leftmargin=*]
  \item \texttt{Gerente}: \texttt{cédulaProfesional} y la relación
        \texttt{Administrar};
  \item \texttt{Cajero}: \texttt{horarioTrabajo}, \texttt{Vender} y
        \texttt{Reflejar};
  \item \texttt{Encargado de premios}: \texttt{horarioTrabajo} y
        \texttt{Procesar};
  \item \texttt{Técnico}: \texttt{áreaEspecialidad}, la certificación
        multivalorada y \texttt{Realizar}.
\end{itemize}

\texttt{Tipo de juego} no es una especialización de \texttt{Juego}; el
diagrama lo modela como una entidad independiente vinculada mediante
\texttt{Clasificar}. Esto permite conservar nombre, regla de puntuación y
condición como datos configurables, sin crear una subentidad por modalidad.
